% Documentation for jigsaw-generator
% Copyright (C) 2015 Julian Gilbey
% This is free software, and you are welcome to redistribute it
% under certain conditions; see the COPYING file for details.

\documentclass{article}
\usepackage{graphicx}
\usepackage{hyperref}
\usepackage{calc}
\usepackage{tikz}[2013/12/13 v3.0.0]
\usetikzlibrary{calc}
\usepackage{amsmath}
\usepackage{upquote}  % to get straight quote marks for code examples

\def\YAML{{\small YAML}}
\def\CMEP{{\small CMEP}}
\def\JSON{{\small JSON}}
\def\PDF{{\small PDF}}
\def\HTML{{\small HTML}}
\def\XML{{\small XML}}
\def\MathML{Math{\small ML}}
\def\MacOSX{Mac\,{\small OS\,X}}
\def\GNU{{\small GNU}}
\def\WYSIWYG{{\small WYSIWYG}}
\def\PGF{{\small PGF}}
\def\TikZ{Ti\emph{k}Z}
\def\CSS{{\small CSS}}
\def\CTAN{{\small CTAN}}
\def\TUG{{\small TUG}}

\usepackage{shortvrb}
\MakeShortVerb{\|}

\def\gobble#1{}

%%%%%%%%%%%%%% The following are macros required for typesetting the
%%%%%%%%%%%%%% jigsaw example.  They are reduced from the full macros
%%%%%%%%%%%%%% which come with the jigsaw software.
\tikzset{regular/.style={sloped, above}}

\newlength\sidelength % side length
\newlength\trad % triangle circumradius
\newcommand{\setshapesize}[1]{%
 \pgfmathsetlength\sidelength{#1*72.27/2.54}
 \pgfmathsetlength\trad{(#1*72.27/2.54) / sqrt(3)}
}

\newcommand{\makestyles}[3]{%
    % split #3 into {style}{text}
    \makestylesa#1#2#3
}
\newcommand{\makestylesa}[4]{%
  \def#1{#3}
  \def#2{#4}
}

\newcommand{\tiltedtriangle}[7]{%
  \coordinate (a) at ($ #1 +({#2-150}:\trad) $);
  \coordinate (b) at ($ #1 +({#2-30}:\trad) $);
  \coordinate (c) at ($ #1 +({#2+90}:\trad) $);

  \makestyles\stylea\contenta{#3}
  \makestyles\styleb\contentb{#4}
  \makestyles\stylec\contentc{#5}

  \draw[allow upside down]
    (a) -- node[\stylea] {\contenta}
    (b) -- node[\styleb] {\contentb}
    (c) -- node[\stylec] {\contentc}
    cycle;
  % If card number is empty, don't draw circle around number
  \def\templabel{#6}
  \ifx\templabel\empty
  \else
    \node[circle,draw,thin,inner sep=1pt,rotate=#7] at ($ #1 $) {#6};
  \fi
}

\newcommand{\uptriangle}[6]{%
  \tiltedtriangle{#1}{0}{#2}{#3}{#4}{#5}{#6}
}
\newcommand{\downtriangle}[6]{%
  \tiltedtriangle{#1}{180}{#2}{#3}{#4}{#5}{#6}
}
\newcommand{\lefttriangle}[6]{%
  \tiltedtriangle{#1}{90}{#2}{#3}{#4}{#5}{#6}
}
\newcommand{\righttriangle}[6]{%
  \tiltedtriangle{#1}{270}{#2}{#3}{#4}{#5}{#6}
}

%%%%%%%%%%%%%% end jigsaw macros


\begin{document}

\title{Creating (mathematical) jigsaw puzzles using jigsaw-generator}
\author{Julian Gilbey}
\date{December 2015}
\maketitle

\begin{abstract}
  \noindent
  The jigsaw puzzles described here are constructed from shapes such
  as triangles, squares and so on, with questions and answers written
  along their edges.  The aim is to match them up correctly.  Related
  puzzle varieties are card sorts and dominoes.  This package provides
  a \TeX- and Python-based system designed to author such puzzles.
  The input is a text-based \YAML\ file; the output can include both
  printable \PDF s for cutting up and Markdown files for potential
  conversion to \HTML.
\end{abstract}

\tableofcontents

\section{Introduction}

For a general background to this software and its history, see
my \TUG boat article~\cite{JGtugboat}.

This documentation describes the usage of the software; some details
have changed since the \TUG boat article, and so all of the relevant
material will be repeated here.

\subsection{Usage overview}

The most important files from a user's perspective are the data input
files.  These specify the type of puzzle (a hexagonal jigsaw or a card
sort, for example), any options for the jigsaw (such as whether to
number the cards or not) and the text to appear on the puzzle.
Sensible defaults are provided for all of the options if they are not
specified.  These input files are written in \YAML.

For example, here is a data file for creating a small hexagonal puzzle:

\begin{verbatim}
format: 1
type: smallhexagon
title: An example puzzle
note: 'You will have to work out the missing number shown as `?'''

pairs:
  - ['$10^6\div10^3$', '$1000$']
  - ['$10^{23}\div10^{21}$', '$100$']
  - ['$9^4\times9^5$', '$9^9$']
  - - '$(3^2)^2$'
    - puzzletext: '?'
      solutiontext: '$81$'
  - ['$6\times6^2$', '$216$']
  - ['$2^4\div2^6$', '$\dfrac{1}{4}$']
edges:
  - '$-1$'
  - '$10$'
  - '$13^2\times13$'
  - '$39$'
  - '$13^3$'
  - '$108$'
\end{verbatim}

The solution to the resulting puzzle looks like this:

\begin{center}
  \tikzset{above/.default=2pt}
  
  \setshapesize{3.7}
  \begin{tikzpicture}[line width=1.5pt, rounded corners=0.7pt]
    \righttriangle{(60:\trad)}  {{regular}{\normalsize $1000$}}{{regular}{\normalsize $10^{23}\div10^{21}$}}{{regular}{\normalsize $108$}}{\scriptsize 1}{150} % {A1}{Q2}{E1}{1}{270}
    \lefttriangle {(0:\trad)}   {{regular}{\normalsize $13^3$}}{{regular}{\normalsize $100$}}{{regular}{\normalsize $9^4\times9^5$}}{\scriptsize 3}{-150} % {E2}{A2}{Q3}{2}{90}
    \righttriangle{(-60:\trad)} {{regular}{\normalsize $(3^2)^2$}}{{regular}{\normalsize $39$}}{{regular}{\normalsize $9^9$}}{\scriptsize 2}{150} % {Q4}{E3}{A3}{3}{270}
    \lefttriangle {(-120:\trad)}{{regular}{\normalsize $81$}}{{regular}{\normalsize $6\times6^2$}}{{regular}{\normalsize $13^2\times13$}}{\scriptsize \underline{6}}{-30} % {A4}{Q5}{E4}{4}{90}
    \righttriangle{(180:\trad)} {{regular}{\normalsize $10$}}{{regular}{\normalsize $216$}}{{regular}{\normalsize $2^4\div2^6$}}{\scriptsize 5}{150} % {E5}{A5}{Q6}{5}{270}
    \lefttriangle {(120:\trad)} {{regular}{\normalsize $10^6\div10^3$}}{{regular}{\normalsize $-1$}}{{regular}{\normalsize $\frac{1}{4}$}}{\scriptsize 4}{-150} % {Q1}{E6}{A6}{6}{90}
  \end{tikzpicture}
\end{center}

If the puzzle is stored in the file |hexpuzzle.yaml|, then the
command:
\begin{verbatim}
jigsaw-generate hexpuzzle
\end{verbatim}
will produce several \PDF\ files: |hexpuzzle-puzzle.pdf|
(containing the puzzle), |hexpuzzle-solution.pdf| (the solution)
and |hexpuzzle-table.pdf| (with the puzzle content laid out in a
tabular form), along with Markdown versions |hexpuzzle-puzzle.md|
and |hexpuzzle-solution.md| containing the puzzle and solution
respectively; these can be used to create \HTML\ versions of the
puzzle via pandoc or some such system.

The command-line command |jigsaw-generate| offers a number of command
line switches; run the command
\begin{verbatim}
jigsaw-generate --help
\end{verbatim}
to obtain information about this.

\bigskip

In the following sections, the data files are described in detail, and
then the other supporting files (configuration files, templates and so
on) and command line options are described.

In many cases, it is adequate to simply copy and edit the example
input files, which will be found in
\nolinkurl{/usr/local/share/doc/jigsaw-generator/examples} or
\nolinkurl{/usr/share/doc/jigsaw-generator/examples}, depending on
your specific installation of the package.

\section{Data input files}

The data input files consist of layout information (the type of
puzzle, whether to produce a solution and so forth) and content
information.  Being a \YAML\ mapping, these can appear in any order in
the data file; will will describe them in these two separate sections
for simplicity.

\subsection{Available layouts}

There are four jigsaw puzzle layouts currently available:
\begin{itemize}
\item smallhexagon, which consists of 6 triangles as shown above;
\item hexagon, which consists of 24 triangles;
\item triangle, which consists of 16 smaller triangles;
\item parquet, which consists of 4 squares and 8 triangles.
\end{itemize}

There are also three card sort layouts currently available:
\begin{itemize}
\item cards, which consists of an arbitrary number of rectangular
  cards; no solution is produced, and this is intended for
  classification activities or suchlike, where there is no ``correct''
  answer expected; these cards are not shuffled by default;
\item cardsort, which is like cards, except that there is a correct
  order, so the cards are shuffled and a solution is produced;
\item dominoes, which consists of an arbitrary number of domino cards;
  a solution is produced.
\end{itemize}

More details are given in the following sections.

\subsection{Layout information}

Each of these options must be followed by a colon and a space,
followed by the option value, as demonstrated in the example above.
All of these options are optional unless stated otherwise.  The most
common options are given first in each of the following three lists.

\subsubsection*{Options applicable to all puzzle types}

\begin{itemize}
\item |format:| The version of the file format.  It is currently
  ignored, but the format version of the current software is 1.
  Future changes to the software will use the format version where
  relevant to produce backward-compatible behaviour for older format
  files.
\item |type:| (required) The type of jigsaw being created.  The
  current recognised standard types are those listed above, namely
  |smallhexagon|, and so on.
\item |title:| The title of the activity, printed at
  the top of the page.
\item |note:| A note printed at the top of the page.
\item |produceSolution:| Whether to produce a solution \PDF\ and
  Markdown file; the default is |true| except for cards, where
  the default is |false|.
\item |numberCards:| Do we number the cards for ease of
  identification?  (Default: |true|)
\item |textSize:| How large the text should be on the cards; this
  is a number from 0 to 9, with 0 being the smallest.  (Default:~5)
\item |makepdf:| Whether to produce \PDF\ output files.  (Default:
  |True|)
\item |makemd:| Whether to produce Markdown output files.  (Default:
  |True|)
\item |latex:| Which \LaTeX\ engine to use to produce the \PDF\
  files.  (Default: |pdflatex|)
\item |texfilter:| If specified, this is a filter to pass the \LaTeX\
  files through prior to running \LaTeX.  See the later section on
  filters (\ref{sec:filters}) for more information on filters.
\item |mdfilter:| If specified, this is a filter to pass the Markdown
  files through to produce the final Markdown files.
\item |clean:| If true, then intermediate files will be deleted.
  This includes the \LaTeX\ files and related aux files and so on.
  (Default: |true|)
\item |puzzleHeaderTeX:| This specifies the header file to use to
  produce the \LaTeX\ puzzle file.  See the later section on templates
  (\ref{sec:templates}) for more information on header and template
  files.
\item |puzzleTemplateTeX:| This specifies the template file to
  use to produce the \LaTeX\ puzzle file.
\item |solutionHeaderTeX:| This specifies the header file to
  use to produce the \LaTeX\ solution file.
\item |solutionTemplateTeX:| This specifies the template file to
  use to produce the \LaTeX\ solution file.
\item |tableHeaderTeX:| This specifies the header file to
  use to produce the \LaTeX\ table file.
\item |tableTemplateTeX:| This specifies the template file to
  use to produce the \LaTeX\ table file.
\item |puzzleHeaderMarkdown:| This specifies the header file to
  use to produce the Markdown puzzle file.
\item |puzzleTemplateMarkdown:| This specifies the template file to
  use to produce the Markdown puzzle file.
\item |solutionHeaderMarkdown:| This specifies the header file to
  use to produce the Markdown solution file.
\item |solutionTemplateMarkdown:| This specifies the template file to
  use to produce the Markdown solution file.
\item |hiddennote:| If any entries are hidden, the note to include at
  the top of the \PDF\ solution output file.  (Default:
  |'|\texttt{Entries that are hidden in the puzzle are highlighted in
    yellow.}|'|)
\item |hiddennotemd:| If any entries are hidden, the note to include
  at the top of the Markdown solution output file.  (Default:
  |'|\texttt{Entries that are hidden in the puzzle are indicated with
    (*).}|'|)
\item |hiddennotetable:| If any entries are hidden, the note to
  include at the top of the \PDF\ table output file.  (Default:
  |'|\texttt{Entries that are hidden in the puzzle are indicated with
    (*).}|'|)
\end{itemize}

\subsubsection*{Options applicable to card sorts and dominoes only}

\begin{itemize}
\item |rows:| (required) How many rows of cards to produce per
  page.
\item |columns:| (required) How many columns of cards to produce
  per page.
\item |shuffleCards:| (card sorts only) Are the cards shuffled?
  This is by default |true| for cardsort, and |false| for
  cards.
\item |shufflePairs:| (dominoes only) Are the order of the pairs
  shuffled from the order given in the data file?  This is by default
  |false|.
\item |flip:| (dominoes only) Are the order of the entries in
  each pair presented in a random order?  This is by default
  |false|.
\item |cardTitle:| The title to write on each card.
\item |cardTitleSize:| How large the title should be on the
  cards, a number from 0 to 9 (like the |textSize|).  By default,
  this is 1~less than the label size (see below).
\item |label:| The default label to write on each card.  This is
  like a subtitle, and appears at the bottom of the card.  (This can
  be useful if the cards are of two or more types, and the type of
  card is important.)  This is not used for dominoes.
\item |labelSize:| How large the labels should be on the cards,
  again a number from 0 to 9; by default, this is 2~less than the text
  size.  (This can be overridden for individual cards.)  This is not
  used for dominoes (though setting it will have an impact on the size
  of the title if |cardTitleSize| is not set).
\item |cardsep:| How much space to leave between
  adjacent cards.  By default this is 12pt.  It should be a \TeX\
  length.
\item |cardsepHorizontal:| How much space to leave between
  horizontally adjacent cards.  By default, this is |cardsep|.
\item |cardsepVertical:| How much space to leave between
  vertically adjacent cards.  By default, this is |cardsep|.
\item |loop:| (dominoes only) Do we have a domino loop, with
  there being no start and end?  The default is |true|.
\item |start:| (dominoes only) If we do not have a loop, then the
  text to write on the start of the chain.  The default is |'Start'|.
\item |finish:| (dominoes only) If we do not have a loop, then the
  text to write on the end of the chain.  The default is |'Finish'|.
% \item \textbf{WILL WANT SOMETHING HERE ABOUT \% DIVIDER ON DOMINOES}
\end{itemize}

\subsubsection*{Options applicable to jigsaws only}

\begin{itemize}
\item |shufflePairs:| Are the order of the pairs shuffled from
  the order given in the data file?  This is by default |false|.
\item |flip:| Are the order of the entries in each pair presented
  in a random order?  This is by default |true|.
\item |shuffleEdges:| Are the order of the edges shuffled from
  the order given in the data file?  This is by default |false|.
\end{itemize}

\subsection{Card content}

The content of the cards is given in a final one or two \YAML\
key--value pairs, depending on the type of material being produced.
The keys are |cards| (for card sorts), |pairs| (for dominoes and
jigsaws) and |edges| (for jigsaws).

Each of these takes as a value a \YAML\ sequence of card content.  For
|cards| and |edges|, each of these is a single entry, while for
|pairs|, each one is a sequence of two entries, as in the example
above.  The number of cards, edges or pairs depends upon the type of
puzzle being produced.

Each entry can have exactly one of the following types of content:

\begin{itemize}
\item A string.  This will be the content of the relevant card, edge
  or half of a pair.  This will be inserted verbatim into the Markdown
  and \LaTeX\ output files, so should be written in \LaTeX.  (Images
  are handled specially; see below.)
\item A \YAML\ mapping.  This uses the following keys:
  \begin{itemize}
  \item |text:| This contains the text to be used as content for this
    entry (card, edge or half of a pair).
  \item |size:| This will be added to the |defaultsize|, so a size of
    |1| will produce larger text in the \PDF\, while a size of |-1|
    will produce smaller text.
  \item |hidden:| If this is set to |true|, the entry will not be
    shown on the puzzle, only on the solution.
  \item |label:| This overrides the current label.
  \item |puzzletext:|, |puzzlesize:| These override the |text| and
    |size| settings for the puzzle.  Specifying |puzzletext| overrides
    any |hidden| setting.
  \item |solutiontext:|, |solutionsize:| These override the |text| and
    |size| settings for the solution.  If |solutiontext| is specified,
    then |hidden| is assumed to be |true|.
  \end{itemize}
  Either |text| or both |puzzletext| and |solutiontext| must be
  present.

\item Special \YAML\ content.
  \begin{itemize}
  \item |newpage: true| This produces a new page in the \PDF\ output
    file for card sort activities, but has no effect on the table or
    Markdown outputs.
  \item |newlabel:| This contains the text to be used as the default
    label for all following cards.
  \item |newlabelsize:| This specifies the label size to be used as
    the default for all following cards.
  \end{itemize}
\end{itemize}



\section{Puzzle templates}
\label{sec:templates}

The puzzle templates are \LaTeX\ files with templating marks.  For the
supplied templates, \PGF/\TikZ\ has been used for the graphics and to
place the text items.  Version 3.0.0 or higher of \PGF/\TikZ\ is
required for the current templates.  Alternative templates can be
written if this is desired.  For example, someone might prefer to use
Asymptote or another graphical package, or they may wish to modify the
existing templates in various ways.

There is also a layout file (written in \YAML) for each puzzle type:
this specifies various parameters required to create the puzzle.

\textbf{TODO: Write a description of the template and layout files}

In the meantime, please look at the template and layout files in the
distributed package to see how they are written, and modify them
appropriately if you wish to create your own.  You can also email me
if you have questions.

These files can be specified in the puzzle files as described above
(with the template description file specified by the |type| option,
and the templates defined within the layout file or overridden by the
corresponding options in the puzzle file).

These files are looked for in the current directory, in the user config
directory (typically |~/.config/jigsaw-generator/templates/|) and in the
package templates directory (typically
|/usr/share/jigsaw-generator/templates/| or
|/usr/local/share/jigsaw-generator/templates/|).

\section{Filters}
\label{sec:filters}

There may be occasions on which you need to systematically modify the
\LaTeX\ file prior to running \LaTeX, or to modify the Markdown output
in some way.  Filters offer this facility.  They are programs (scripts
or other executables) which read in the \LaTeX\ or Markdown on
standard input and write the filtered version to standard output.  For
example, you might have reason to change the standard Markdown or
\LaTeX\ syntax for images to meet some local need.

Filters can been specified either on the command line or within the
puzzle files or within the template description files.

Filters are looked for in the current directory, in the user config
directory (typically |~/.config/jigsaw-generator/filters/|) and in the
package filters directory (typically
|/usr/share/jigsaw-generator/filters/| or
|/usr/local/share/jigsaw-generator/filters/|; this assumes that all
supplied filters are system-independent scripts).

\section{Configuration files}

The |jigsaw-generate| program reads a configuration file,
|~/.config/jigsaw-generator/config.ini|, which is in an INI format.
The default configuration file can be found in
|/usr/share/jigsaw-generator/templates/default_config.ini| (or
|/usr/local/share/...|, depending on where is has been installed).

The default file is effectively as follows:

\begin{verbatim}
[jigsaw-generate]
makepdf = yes
makemd = yes
latex = pdflatex
clean = yes
texfilter = 
mdfilter = 
\end{verbatim}

and these correspond to the identically-named command-line options.

\section{Some notes on \texorpdfstring{\YAML}{YAML} syntax}

For a full, precise description of \YAML\ syntax, see the \YAML\
Specification~\cite{YAML}.  What follows is a very brief summary of
some of the basic parts of the specification which should be
sufficient for most people's needs when using this software.

Although \YAML\ allows unquoted strings in general, there are a number
of restrictions on what is permitted in them.  For this reason, it may
be simplest to single-quote all value strings.  (\YAML\ also offers
double-quoted strings, but these interpolate backslash-escapes; this
is probably undesirable in this context, since many \TeX\ expressions
include backslashes.)

Within a single-quoted string, a single quote is written as a doubled
quote mark (|''|), hence the three quote marks in a row at the
end of the note field in the above example (two to indicate a quote,
and the third to end the string).

For collections (`sequences' and `mappings' in \YAML's terminology,
each of which consists of a number of `entries'), \YAML\ allows two
different notations: either a flow style, which is similar to \JSON\
notation, or a block style.  With the flow style, the entries of a
sequence, such as a question--answer pair or the list of edges,
are enclosed in square brackets and separated by commas; for the block
style, each entry appears on a new line preceded by a
vertically-aligned hyphen.

Similarly, for a mapping the entries consist of key--value pairs, with
the key and value separated by a colon.  They can be written either
using a flow style as a comma-separated list enclosed in braces, or
with a block style by writing each key--value pair on an
identically-indented new line.

Both of the sequence styles appear in the example here, though only
the block style is used for mappings.  Note also that the top-level
structure of the file is itself a mapping.  This means that the
entries (type, title, pairs, and so on) can appear in any order.
However, for the benefit of the human reader, it is wise to maintain a
meaningful order to these entries.

\section*{Acknowledgements}

I thank those who offered very useful ideas and feedback at the \TUG\
conference; these have helped me to solve some of the thorny issues I
had been facing.

\smallskip
\bibliographystyle{plain}
\bibliography{jigsawrefs}

\end{document}

%%% Local Variables: 
%%% mode: latex
%%% TeX-master: t
%%% End: 
